% ----------------------------------------------------------
% Trabalhos Futuros
% ----------------------------------------------------------
\chapter{Trabalhos Futuros}

Portanto, sobre o que o sistema tem a oferecer, objetiva-se elencar as melhorias e avanços necessários para otimizar a ferramenta, contribuindo para uma experiência mais completa e segura ao usuário. Dentre as melhorias previstas para o sistema, destacam-se:

\begin{itemize}
    \item Implementação da extração automática de dados a partir de documentos oficiais do processo, permitindo o autopreenchimento  do formulário de execução;
    \item Implementação da funcionalidade de unificação de penas, permitindo que condenações distintas sejam consolidadas em um único cálculo de execução;
    \item Desenvolvimento de controle de acesso por perfil de usuário, diferenciando as permissões entre advogados, defensores e demais operadores do sistema;
    \item Implementação da verificação de e-mail no cadastro, garantindo maior segurança no acesso à plataforma;
    \item Integração com sistemas oficiais, permitindo a sincronização direta com os processos já cadastrados nessas plataformas; e
    \item Implementação do direito de exclusão de dados pelo titular 
    e elaboração de uma política de privacidade acessível na plataforma, em conformidade com a LGPD \cite{lgpd}.
\end{itemize}

Portanto, as melhorias destacadas têm como objetivo aprimorar o sistema CalPEC, com a finalidade de abranger de forma ainda mais eficiente de extender os usuários finais: advogados, defensores públicos e demais operadores 
do Direito. Assim, o sistema reforça seu diferencial em relação às demais ferramentas existentes, transmitindo utilidade e segurançado do sistema no cotidiano forense.